\subsection{Image--Label Fusion}\label{sec:results-fusion}

This section compares three architectures. The comparison between Dual U-Net
 and shared-encoder architectures tests the change from two U-Nets coupled by
cross-attention to a shared encoder with attention layers. The comparison
between single axis and dual axis attention tests whether, within the
shared-encoder design, keeping image and label tokens separate improves over
merging them once.

\paragraph{Architectures.}
\emph{Dual U-Net Cross-Attn} is Medverse's architecture
(\autoref{sec:fundamentals-relatedwork}), initialized from the released weights
and then trained with our recipe. It is therefore a different model from the
released Medverse baseline used in \autoref{sec:results-cascade}.
\emph{Single axis} and \emph{Dual axis attention} are the two variants shown in
\autoref{fig:fusion-axis}. Early fusion uses $N_\text{layer}{=}9$ layers to match the compute of dual attention (\autoref{sec:methodology-fusion}).
This gives both variants similar compute (3973 vs.\ 3892 GFLOPs). Parameter
count could not be matched at the same time, and early fusion has 52\% more
parameters. Dual U-Net Cross-Attn is not compute-matched to either variant.

\paragraph{Training.}
All three architectures use the same training recipe
(\autoref{tab:fusion-trainsettings}), data pipeline, and augmentations. They
are trained on the CT cases and training classes of TotalSegmentator, at a single spacing level randomly sampled in the range 1.5--6\,mm. 
Dual U-Net
Cross-Attn starts from Medverse's weights. Early fusion and dual attention
are randomly initialized.

\begin{table}[t]
  \centering
  \small
  \caption{Training settings shared by the three architectures of
  \autoref{sec:results-fusion}.}
  \label{tab:fusion-trainsettings}
  \begin{tabular}{@{}ll@{}}
    \toprule
    Optimizer & AdamW \\
    Learning rate & $10^{-4}$ \\
    Schedule & cosine decay, 10-epoch warmup \\
    Weight decay & $10^{-2}$ \\
    Batch size & 8 \\
    Gradient clipping & 1.0 \\
    Loss & BCE + Dice \\
    Epochs & 250 \\
    \bottomrule
  \end{tabular}
\end{table}

\paragraph{Evaluation.}
All three architectures are evaluated at a single level with 3\,mm spacing,
with crops centered on the ground truth, so localization is not tested here.
\autoref{tab:fusion-staircase} reports the results.


\begin{table}[t]
  \centering
  \small
  \caption{Dice of the three architectures at a single level (3\,mm), with
  parameter count, compute, and latency.}
  \label{tab:fusion-staircase}
  \begin{tabular}{@{}lccc@{}}
    \toprule
    & Dual U-Net & Single axis Attn & Dual Axis Attn \\
    \midrule
    TotalSeg CT, seen      & 0.348 & 0.574 & \textbf{0.582} \\
    TotalSeg CT, unseen    & 0.221 & 0.408 & \textbf{0.426} \\
    FLARE22, seen          & 0.532 & 0.693 & \textbf{0.709} \\
    FLARE22, unseen        & 0.443 & \textbf{0.632} & 0.604 \\
    \midrule
    MSD Prostate           & \textbf{0.377} & 0.205 & 0.247 \\
    MSD Hippocampus        & 0.172 & 0.122 & \textbf{0.184} \\
    \midrule
    HU\_LWK1               & \textbf{0.146} & 0.079 & 0.102 \\
    \midrule
    TotalSeg MRI           & 0.178 & 0.232 & \textbf{0.259} \\
    \midrule
    ISLES22                & \textbf{0.046} & 0.037 & 0.043 \\
    Shifts-MS               & 0.087 & 0.038 & \textbf{0.109} \\
    ATLAS v2.0              & \textbf{0.031} & 0.014 & 0.022 \\
    GNC\_705                & \textbf{0.044} & 0.014 & 0.025 \\
    NasalSeg                & \textbf{0.434} & 0.272 & 0.358 \\
    \midrule
    Mean of 10 held-out sources & 0.187 & 0.174 & \textbf{0.201} \\
    \midrule
    Parameters             & 71.1M  & 74.1M  & 48.7M  \\
    GFLOPs                 & 2362.6 & 3973.1 & 3892.5 \\
    Latency per sample     & 73.2\,ms & 73.0\,ms & 72.5\,ms \\
    \bottomrule
  \end{tabular}

\end{table}

\paragraph{Results.}
Replacing the two U-Nets with a shared encoder raises Dice by 0.16 to 0.23 on
the seen and unseen classes of TotalSegmentator and FLARE22, and from 0.178 to
0.232 on the unseen modality TotalSeg MRI. On all eight datasets with targets
outside the training vocabulary, however, single-axis attention is below Dual
U-Net Cross-Attn.

Dual attention improves over single-axis attention on 12 of the 13 rows; the
exception is the two unseen FLARE22 classes ($-0.028$). The gains are small on
the in-vocabulary rows (at most 0.018) and largest on new targets: NasalSeg
($+0.086$), Shifts-MS ($+0.071$), MSD Hippocampus ($+0.062$), and MSD Prostate
($+0.042$).

Compared with Dual U-Net Cross-Attn, dual attention is better on the four
in-vocabulary rows, TotalSeg MRI, MSD Hippocampus, and Shifts-MS, and worse on
the other six new-target datasets, most clearly MSD Prostate ($-0.130$) and
NasalSeg ($-0.076$). Its higher held-out mean (0.201 vs.\ 0.187) is thus
driven by FLARE22 and TotalSeg MRI, whose classes are largely in the training
vocabulary; on the remaining eight datasets, Dual U-Net Cross-Attn leads
(unweighted mean 0.167 vs.\ 0.136). Part of this lead may come from its
initialization with Medverse's weights.

Dual attention has the fewest parameters (48.7M); its compute is similar to
single-axis attention by construction and 65\% higher than Dual U-Net
Cross-Attn. Latency is nearly identical for all three (72.5 to 73.2\,ms).